\section*{الفصل \ref{ch:b3:structural-analysis} --- التحليل البنيوي}

\begin{solution}{exo:b3:structural-analysis:1}
يتّجه \omterm{def:b3:structural-analysis:stereo}{القطب} نحو $40 + 180 = \ang{220}$ ويغوص بمقدار $90 - 30 = \ang{60}$. أمّا
المستوي فينحدر بمقدار $90 - 20 = \ang{70}$ نحو $315 - 180 = \ang{135}$: $135/70$.
\end{solution}

\begin{solution}{exo:b3:structural-analysis:2}
$\theta = \ang{60}$. في \omterm{def:b3:structural-analysis:stereo}{الشبكة متساوية الزوايا}: $10\tan\ang{30} = \qty{5.77}{cm}$؛
وفي الإسقاط متساوي المساحات: $10\sqrt2\sin\ang{30} = \qty{7.07}{cm}$.
\end{solution}

\begin{solution}{exo:b3:structural-analysis:3}
$\vec n_1 = (\sin\ang{30}, 0, \cos\ang{30})$، $\vec n_2 = (-\sin\ang{30}, 0,
\cos\ang{30})$، $\vec n_1\cdot\vec n_2 = -0.25 + 0.75 = 0.5$: القطبان متباعدان بمقدار
\ang{60}، وزاوية الأجنحة \ang{120}، وهي طية مفتوحة.
\end{solution}

\begin{solution}{exo:b3:structural-analysis:4}
$\vec d = (\cos\ang{60}, 0, -\sin\ang{60}) = (0.5, 0, -0.866)$ و$\vec u =
0.707\,\vec s + 0.707\,\vec d = (0.354, 0.707, -0.612)$: يتّجه
$\operatorname{atan2}(0.354, 0.707) = \ang{26.6}$، ويغوص بمقدار $\arcsin0.612 =
\ang{37.8}$. تقع زاوية الانزلاق بين \ang{-90} (عادي) و0 (أيسر): صدعٌ
عاديّ مائل مع مركّبة أيسر.
\end{solution}

\begin{solution}{exo:b3:structural-analysis:5}
\begin{gather*}
  \vec n_1 = (0.5567, 0.3214, 0.7660),\qquad \vec n_2 = (-0.4698, -0.1710, 0.8660),\\
  \vec n_1\times\vec n_2 = (0.4093, -0.8420, 0.0558),
\end{gather*}
بطول $0.9379$. ومعكوسه، $-\vec n_1\times\vec n_2$، يشير إلى الأسفل:
يتّجه نحو \ang{334.1}، ويغوص بمقدار $\arcsin(0.0558/0.9379) = \ang{3.4}$. $\vec n_1\cdot\vec
n_2 = 0.3469$، و\ang{69.7} بين العموديين، وزاوية الأجنحة \ang{110.3}:
طية مفتوحة.
\end{solution}

\begin{solution}{exo:b3:structural-analysis:6}
الغوص فوق \ang{60} يعني $\theta < \ang{30}$. في الإسقاط متساوي المساحات: $r^2 =
2\sin^2\ang{15} = 0.134$، \qty{13.4}{\percent} من الشبكة؛ وفي \omterm{def:b3:structural-analysis:stereo}{الشبكة متساوية الزوايا}:
$\tan^2\ang{15} = 0.072$، \qty{7.2}{\percent}. على الكرة تحتوي القلنسوة
$1 - \cos\ang{30} = 0.134$ من نصف الكرة: لا تُظهرها إلا الشبكة متساوية المساحات
بحصّتها الحقيقية، ولذلك تُعدّ الأقطاب عليها.
\end{solution}

\begin{solution}{exo:b3:structural-analysis:7}
\begin{gather*}
  T_{xx} = \frac{2\sin^2\ang{20}}3 = 0.078,\qquad
  T_{zz} = \frac{2\cos^2\ang{20} + 1}3 = 0.922,\\
  T_{xz} = -\frac{\sin\ang{20}\cos\ang{20} - \sin\ang{20}\cos\ang{20}}3 = 0,
\end{gather*}
وكل حدّ يتضمّن $y$ يساوي 0. القيم الذاتية هي $0.922$ (رأسية)،
$0.078$ (شرق) و$0$ (شمال): \omterm{def:b3:structural-analysis:pi}{محور بيتا} أفقي، شمال--جنوب.
تقع الأقطاب الثلاثة في مستوٍ واحد، فلا يكون لأيّ منها مركّبة على عموديه،
و$\vec u^{\mathsf T}T\vec u = 0$ من أجل $\vec u$ الموازية له.
\end{solution}

\begin{solution}{exo:b3:structural-analysis:8}
$a = 2$، $b = 1.5$، $k = 1/0.5 = 2$: تضيّق. $Z = 1/(1.8\times1.1) =
0.505$، $a = 1.636$، $b = 2.178$، $k = 0.636/1.178 = 0.54$: التفلطح.
\end{solution}

\begin{solution}{exo:b3:structural-analysis:9}
$S_1/S_2 = 2.25$ و$S_1S_2 = 1$: $S_1 = 1.5$، $S_2 = 0.667$؛ المحور الطويل
استطال بمقدار \qty{50}{\percent}. للفراغ أنصاف محاور $1.5\times0.4 =
\qty{0.60}{mm}$ و$0.667\times0.4 = \qty{0.27}{mm}$.
\end{solution}

\begin{solution}{exo:b3:structural-analysis:10}
في \omterm{def:b3:structural-analysis:stereo}{الشبكة متساوية الزوايا}: $1/(2\cos^2\ang{30}) = 0.667$ على طول نصف القطر،
و$\tan\ang{30}/\sin\ang{60} = 0.577/0.866 = 0.667$ في الاتجاه المماسّي: متساويان، والإسقاط مطابِق.
في الإسقاط متساوي المساحات: $\mathrm dr/(R\,\mathrm d\theta) = (\sqrt2/2)\cos\ang{30} = 0.612$
و$\sqrt2\sin\ang{30}/\sin\ang{60} = 0.707/0.866 = 0.816$؛ وحاصل الضرب
$0.500$، وهو معامل المساحة في المبرهنة. تصبح الدائرة الصغرى بيضويًّا، محاوره
بنسبة $0.816/0.612 = 1.33$، والمحور الأطول حول المركز.
\end{solution}

\begin{solution}{exo:b3:structural-analysis:11}
الصدع رأسي، فعموديه أفقي ومتعامد مع الإجهاد
الرأسي $\sigma_2$؛ وحسب الجزء (ب) من المبرهنة يكون الانزلاق عموديًّا على
$\sigma_2$، أفقيًّا، ومن أجل كلّ $\Phi$. مع $\sigma_1 = 1$ باتجاه الشمال،
$\sigma_3 = 0$ باتجاه الشرق: $S\vec n = (0, -0.342, 0)$، $\vec n\cdot S\vec n =
0.117$، $\vec\tau = (0, -0.342, 0) - 0.117\,(0.940, -0.342, 0) = (-0.110,
-0.302, 0)$ و$-\vec\tau = (0.110, 0.302, 0)$، باتجاه \ang{020}. الكتلة
الواقعة على جانب $\vec n$، أي الجانب الشرقي، تتحرّك نحو الشمال الشمالي الشرقي على امتداد
الامتداد: عند النظر من الجانب الغربي تتحرّك نحو اليسار، فهو صدعٌ أيسر (زاوية الانزلاق
0). لم يدخل $\sigma_2$ في الحساب.
\end{solution}

\begin{solution}{exo:b3:structural-analysis:12}
$S = (7.5 - 6.0)/1.0 = \qty{1.5}{km}$؛ و$h = 6.0/1.5 = \qty{4.0}{km}$ تحت
الأفق السفلي. كان طول المقطع $20 + 1.5 = \qty{21.5}{km}$، وتغيّر
بمقدار $-1.5/21.5 = \qty{-7.0}{\percent}$.
\end{solution}

\begin{solution}{pb:b3:structural-analysis:1}
\textbf{1.} يتّجه \omterm{def:b3:structural-analysis:stereo}{القطب} نحو \ang{255} ويغوص بمقدار \ang{77}؛ $\theta =
\ang{13}$ و$r = 10\sqrt2\sin\ang{6.5} = \qty{1.60}{cm}$ عن المركز،
باتجاه \ang{255}.

\textbf{2.} $\vec n = (\sin\ang{13}\sin\ang{75}, \sin\ang{13}\cos\ang{75},
\cos\ang{13}) = (0.2173, 0.0582, 0.9744)$. ستّ متجهات وحدة شبه متوازية
يبلغ مجموعها قرابة 6: $|\vec S_E| = 5.9974$.

\textbf{3.} الجناح الشرقي: $\bar\alpha = \operatorname{atan2}(1.2583, 0.2017) =
\ang{80.9}$ و$\bar\delta = \arccos(5.8604/5.9974) = \ang{12.3}$؛ $\bar R
= 5.9974/6 = 0.99956$ والتشتّت هو $\sqrt{2\times0.00044} =
\qty{0.0297}{rad}$، أو \ang{1.7}. الجناح الغربي: $\bar\alpha =
\operatorname{atan2}(-1.2445, 0.1886) = \ang{278.6}$ و$\bar\delta =
\arccos(5.8636/5.9972) = \ang{12.1}$؛ $\bar R = 0.99953$، والتشتّت \ang{1.8}.

\textbf{4.}
\[
  \vec S_E\times\vec S_W = (0.0774, -14.6714, 0.4883),
\]
بطول $14.680$. معكوسه يشير إلى الأسفل: الاتجاه هو $\operatorname{atan2}(-0.0774, 14.6714)$، أي $\ang{359.7}$،
والغوص هو $\arcsin(0.4883/14.680)$، أي $\ang{1.9}$.

\textbf{5.} جيب تمام الزاوية بين العموديين المتوسّطين هو
$32.835/35.968 = 0.9129$، أي الجداء القياسي للمجموعين مقسومًا على
طوليهما: \ang{24.1} بين العموديين، وزاوية الأجنحة \ang{155.9}:
طيّة مفتوحة، أقرب إلى اللطيفة.

\textbf{6.}
\begin{gather*}
  \vec m_E = (0.2098, 0.0336, 0.9772),\qquad \vec m_W = (-0.2075, 0.0314, 0.9777),\\
  \vec m_E - \vec m_W = (0.4173, 0.0022, -0.0005):
\end{gather*}
عمودي أفقي متجه نحو \ang{089.7}. المستوي المحوري رأسي، وامتداده
\ang{359.7}، أي تقريبًا شمال--جنوب: الطيّة قائمة.

\textbf{7.} الحزام هو المستوي الذي قطبه هو \omterm{def:b3:structural-analysis:pi}{محور بيتا}: اتجاه الانحدار
\ang{179.7}، والانحدار $90 - 1.9 = \ang{88.1}$، أي تقريبًا المستوي الرأسي
شرق--غرب. كلّ جناح يحتوي محور الطيّة، لذلك فكلّ عمودي عليه يكون عموديًا على المحور
ويقع في المستوي العمودي عليه.

\textbf{8.} $0.04400 + 0.00144 + 0.95456 = 1.00000$. $\vec b^{\mathsf T}T\vec
b = 0.00038$، أصغر قيمة ذاتية: $\vec b$ هو، بالتقريب، المتجه الذاتي
$\vec e_3$، والحزام المارّ بالقطبين المتوسّطين هو
الحزام الأفضل ملاءمةً لجميع القراءات الاثنتي عشرة.

\textbf{9.} $\arcsin\sqrt{0.00038} = \arcsin0.0195 = \ang{1.1}$.

\textbf{10.} تقع الأقطاب على طول الحزام بنحو \ang{12.2} عن الرأسي،
على الجانبين؛ $\sin^2\ang{12.2} = 0.045$، وهو متوسط مربعات
مركّباتها شرق--غرب، وهو $\lambda_2$.

\textbf{11.} نعم: تنحرف الأقطاب عن الحزام بمقدار \ang{1.1} في المتوسط، لا
أكثر من تشتّت \ang{1.7} لقراءات جناح واحد حول متوسّطها.

\textbf{12.} $\ang{1.1}/\sqrt{12} = \ang{0.3}$: الغوص هو \ang{1.9} $\pm$
\ang{0.3}.

\textbf{13.} زاوية انزلاق \ang{-90}: صدعٌ عادي خالص، انزلاقه انحداري. أندرسون:
$\sigma_1$ رأسي، و$\sigma_3$ أفقي على طول اتجاه الانحدار،
\ang{120}--\ang{300}، و$\sigma_2$ أفقي على طول الامتداد،
\ang{030}--\ang{210}.

\textbf{14.} عمودي المستوي $120/60$ يقع في المستوي الرأسي الممتدّ بالاتجاه
\ang{120}، وهو يحتوي $\vec e_1$ و$\vec e_3$: $n_2 = 0$. حسب الجزء (ب) فإنّ
الانزلاق عمودي على $\sigma_2$، أي الامتداد، مهما كانت $\Phi$: هو
انزلاق انحداري. لا يخبرنا الصدع بشيء عن $\Phi$؛ ففي الصيغة يتلاشى كلّ من $n_2$
و$\vec e_2\cdot\vec w$، وبالتالي البسط والمقام يتلاشيان.

\textbf{15.} $\vec n = (\sin\ang{55}\sin\ang{150}, \sin\ang{55}\cos\ang{150},
\cos\ang{55}) = (0.4096, -0.7094, 0.5736)$; $n_1 = 0.574$; $n_2 = 0.5\times
0.4096 + 0.866\times(-0.7094) = -0.410$.

\textbf{16.} $\vec e_1\cdot\vec w = 0.198$ و$\vec e_2\cdot\vec w = 0.5\times
0.771 + 0.866\times0.605 = 0.909$، ومنه $\Phi = -(0.574\times0.198)/(-0.410\times
0.909) = 0.1137/0.3727 = 0.305$.

\textbf{17.} $(0.305 + 0.294 + 0.301 + 0.306)/4 = 0.30$. قُرّبت زوايا الانزلاق
إلى أقرب درجة؛ تغيّر الدرجة الواحدة من زاوية الانزلاق $\Phi$ بنحو 0.02 في هذه
الصدوع، لذا فإنّ نصف الدرجة من التقريب يفسّر التباين.

\textbf{18.} $\sigma_2 - \sigma_3 = 0.3(\sigma_1 - \sigma_3)$: $\sigma_2$ أقرب
إلى $\sigma_3$ منه إلى $\sigma_1$. كلا الإجهادين الأفقيين أدنى بكثير
من الرأسي؛ امتدّت المنطقة أساسًا نحو \ang{120}، وقليلًا
على طول \ang{030}.

\textbf{19.} مع $\vec n = (0.750, -0.433, 0.500)$ والانزلاق $\vec u =
(0.433, -0.250, -0.866)$، فإنّ $(\vec n - \vec u)/\sqrt2$ يغوص بزاوية \ang{75} باتجاه
\ang{300}: P؛ و$(\vec n + \vec u)/\sqrt2$ يغوص بزاوية \ang{15} باتجاه \ang{120}:
T. تقع في المستوي الرأسي نفسه الذي يضم $\sigma_1$ و$\sigma_3$ لكن على بُعد
\ang{15} عنهما، لأنّ الصدع يصنع \ang{30} لا \ang{45} مع
$\sigma_1$.

\textbf{20.} $\tan a = \tan\ang{12.3}\times\cos(\ang{90} - \ang{80.9}) =
0.2180\times0.9874 = 0.2153$: $a = \ang{12.1}$.

\textbf{21.} $L_0 = 8000/\cos\ang{12} = \qty{8178.7}{m}$, $S =
\qty{178.7}{m}$، وهو تغيّر نسبي مقداره $\cos\ang{12} - 1 =
\qty{-2.2}{\percent}$.

\textbf{22.} $4000\tan\ang{12} = \qty{850.2}{m}$; $A = \tfrac12\times8000\times
850.2 = \qty{3.40e6}{m^2}$.

\textbf{23.} $h = \num{3.40e6}/178.7 = \qty{19.0}{km}$. لو وُجد انفصال عند
قاعدة الغطاء لأعطى $h$ بضع مئات من الأمتار ومساحةً زائدة
أصغر بنحو خمسين مرة: فالطيات غير منفصلة هناك.
فإمّا أن يكون الأساس قد انطوى مع الغطاء، أو أن الطيات المقعّرة هبطت دون
المستوى الأصلي، وهو ما يجعل $A$ أصغر؛ والمقطع وحده لا يستطيع أن يحسم ذلك.

\textbf{24.} الإزاحة الأفقية $200/\tan\ang{60} = \qty{115.5}{m}$، والانزلاق $200/\sin\ang{60}
= \qty{230.9}{m}$.

\textbf{25.} $\tan a = \tan\ang{60}\cos\ang{30} = 1.500$: $a = \ang{56.3}$،
والإزاحة الظاهرية هي $200/1.5 = \qty{133.3}{m}$. ينزلق الصدع
إلى أسفل الانحدار باتجاه \ang{120}: والمقطع على طول \ang{120} يُظهر الأطوال الحقيقية
ولذا ينبغي أن يكون هو المرمَّم.

\textbf{26.} الصدع يقطع الطيات، فهو إذن أحدث، ويُرمَّم أولًا،
فيُعاد الجدار المعلّق صعودًا على طول الصدع بمقدار \qty{230.9}{m}؛ ثم
تُفرد الطيات. تُوضع خطوط التثبيت على الأسطح المحورية للطيات المقعّرة، في
الجدار المرتكز، حيث لم يكن للطيات المتناظرة انزلاق بين الطبقات.

\textbf{27.} يغوص محور طيّة فاردن
\textbf{\ang[text-series-to-math=true]{1.9}} نحو الشمال (الاتجاه
\ang{359.7}): \ang{2} المدمجة في المنطقة، المستعادة إلى
\ang{0.3} التي تسمح بها الاثنتا عشرة قراءة.
\end{solution}
